\documentclass[twocolumn, amsthm]{autart}
\usepackage[T1]{fontenc}
\pdfobjcompresslevel=0
\pdfminorversion=4

\usepackage{notation}
\newcommand{\defining}[1]{\textbf{#1}}

\begin{document}

\begin{frontmatter}
\runtitle{Dynamic Search for Dynamic Extension}
\title{The Dynamic Search for the Minimal Dynamic Extension}

\thanks[footnoteinfo]{This paper was not presented at any IFAC 
meeting. Corresponding author R.~S.~D'Souza.}

\author[RSD]{Rollen S. D'Souza}\ead{rollen.dsouza@uwaterloo.ca}

\address[RSD]{Mississauga, ON, Canada}
          
\begin{keyword}                           % Five to ten keywords,  
Cicero; Catiline; orations.               % chosen from the IFAC 
\end{keyword}                             % keyword list or with the 
                                          % help of the Automatica 
                                          % keyword wizard


\begin{abstract}
  For nonlinear control systems, transformation into a linear control system via feedback linearization is one o
  Of course, the method itself only applies to a small class of nonlinear systems, 
\end{abstract}

\end{frontmatter}

\section{Introduction}
Video, patres conscripti, in me omnium vestrum ora atque oculos esse 
conversos, video vos non solunn de vestro ac rei publicae, verum 
etiam, si id depulsum sit, de meo periculo esse sollicitos. Est mihi 
iucunda in malis et grata in dolore vestra erga me voluntas, sed eam, 
per deos inmortales, deponite atque obliti salutis meae de vobis ac 
de vestris liberis cogitate. Mihi si haec condicio consulatus data 
est, ut omnis acerbitates, onunis dolores cruciatusque perferrem, 
feram non solum fortiter, verum etiam lubenter, dum modo meis 
laboribus vobis populoque Romano dignitas salusque pariatur.

\section{Background}

\subsection{Notation}
% Must define the dimension of a graded subalgebra that is simply, finitely generated.
\(\Manifold{M}\)
\(\OpenSet{U}\)
\(\TangentSpace{M}{p_0}\)
\(\CotangentSpace{M}{p_0}\)
\(\SmoothVectorFields{M}\)
\(\SmoothForms{M}\)
\(\VectorSpace{V}_k\)
\(\Algebra{X}_1\)

\subsection{Categories}

\subsection{Dynamic Feedback Linearization}

\subsection{The Derived Flag}

\section{Main Statement}

\subsection{The Integrability Defect Index}

\subsection{The Proposed Algorithm}

\section{The Hitting Graded Subalgebra Problem}
The integrability defect index reduction step --- defect reduction step, in short --- drives the dynamic search for a minimal dynamic extension.
The search algorithm seeks for what this article terms the ``hitting graded subalgebra'' to a list of graded algebras.
The name comes from a related problem known as that ``hitting subset problem'' which has more recently been further extended to the hitting subspace problem.
These problems ask to find the minimal subset (subspace) such that the intersection with every subset (subspace) in the list is non-trivial.
Both problems are known to be NP-Hard.

The defect reduction step asks slightly more than that traditional hitting problem does.
However, to ease demonstration, we first present a solution to the hitting graded subsalgebra problem and then adjust it to the requirements for defect reduction.

\subsection{Problem Statement}
Fix \(k \geq 1.\) 
Let \(\Algebra{X}\) be a simply, finitely-generated graded algebra, and let \(\Algebra{X}_1,\) \(\ldots,\) \(\Algebra{X}_k \subseteq \Algebra{X}\) be a sequence of simply, finitely-generated graded subalgebras.
The hitting graded subalgebra problem is as follows.
%
\begin{prob}[Hitting Graded Subalgebra]
  \label{prob:HSAP}
  A simply, finitely-generated, graded subalgebra \(\Algebra{X}_\star \subseteq \Algebra{X}\) is said to be a \defining{hitting graded subalgebra} to \(\{\Algebra{X}_1,\) \(\ldots,\) \(\Algebra{X}_k\}\) if it is the smallest (in dimension\footnote{Well-defined since we only consider simply, finitely-generated subalgebras as candidate solutions.}) subalgebra that has a non-trivial intersection with every graded subalgebra \(\Algebra{X}_j\) in the sequence.
\end{prob}

When \(X_1 \cap \cdots \cap X_k\) has dimension greater than zero,
any one-dimensional subalgebra of the intersection is a hitting graded subalgebra.
In general, a hitting graded subalgebra \(Y\) solves the problem,
%
\begin{equation}
  \begin{array}{cc}
    {\displaystyle \min_{Y \subseteq X}{}} & \Dimension(Y)\\
    \mathrm{s.t.}
      &  \Dimension(Y \cap X_1) \geq 1\\
      &  \vdots\\
      &  \Dimension(Y \cap X_k) \geq 1\\
  \end{array}
\end{equation}
%
There is no unique solution to this problem in general.
For instance, if \(X_1 \cap \cdots \cap X_k\) has dimension greater than one, then there exists infinitely many solutions.
However, there always exists a solution when all the subalgebras are non-trivial.
%
\begin{lem}
  Problem~\ref{prob:HSAP} has at least one solution if, and only if, \(\Dimension(X_i) \neq 0\) for all \(1 \leq i \leq k.\) \hfill
\end{lem}

Constructing solutions to Problem~\ref{prob:HSAP} is not trivial.
The problem is inherently combinatorial when the intersection of all the subalgebras is trivial.
This can be explicitly shown as Problem~\ref{prob:HSAP} is reducible to the hitting subset problem in a manner akin to the hitting subspace problem.
%
\begin{thm}
  \label{thm:nphard}
  Problem~\ref{prob:HSAP} is NP-Hard.
\end{thm}
%
Theorem~\ref{thm:nphard} suggests the problem can be solved with any practically-efficient algorithm used for NP-Hard problems such as Branch-and-Bound.
Unfortunately, it is not that simple.

Problem~\ref{prob:HSAP} is written in a manner akin to the hitting subset and subspace problems.
\emph{Any} of the minimal size solutions are valid.
Recall that there is usually no unique solution.
However, this problem arises in the context of the defect reduction step which is executed iteratively on the defective ideals;
it is not readily apparent that the choice of which generators to dynamically extend matters or not for future defect reduction steps!
In fact, it is the case that a poor choice amongst the minimal choices of defect reduction leads to an overall non-minimal dynamic extension for the whole system.

In light of this, we now present a dynamic programming approach to generating multiple solutions to Problem~\ref{prob:HSAP}.
Later, Section~\ref{DP-DFL} extends this dynamic algorithm to solve the larger problem of the minimal dynamic extension.

\subsection{The Principle of Optimality}


\subsection{The Dynamic Search for the Hitting Graded Subalgebras}


\section{The Defect Reduction Step}


\section{The Dynamic Search for the Minimal Dynamic Extension}
\label{sec:DP-DFL}


\section{Discussion}


\begin{figure}
\begin{center}
% \includegraphics[height=4cm]{jcaesar.eps}    % The printed column  
\caption{Gaius Julius Caesar, 100--44 B.C.}  % width is 8.4 cm.
\label{fig1}                                 % Size the figures 
\end{center}                                 % accordingly.
\end{figure}

% OR

%\begin{figure}
%\begin{center}
%\epsfig{file=jcaesar,width=7cm}
%\caption{Gaius Julius Caesar, 100--44 B.C.}
%\label{fig1}
%\end{center}
%\end{figure}


\subsection{A subsection}
Marcus Tullius Cicero, 106--43 B.C. was a Roman statesman, orator, 
and philosopher.  A major figure in the last years of the Republic, 
he is best known for his orations against Catiline\footnote{
This footnote should be very brief.}
and for his mastery of Latin prose \cite{Heritage:92}. He was a 
contemporary of Julius Caesar (Fig.~\ref{fig1}).

\section{The argument}
Some words might be appropriate describing equation~(\ref{e1}), if 
we had but time and space enough.
\begin{equation} \label{e1}
{{\partial F}\over {\partial t}} =
D{{\partial^2 F}\over {\partial x^2}}.
\end{equation}
See \cite{Abl:56}, \cite{AbTaRu:54}, \cite{Keo:58} and 
\cite{Pow:85}.
This equation goes far beyond the celebrated theorem ascribed to the great
Pythagoras by his followers.
\begin{thm}
The square of the length of the hypotenuse of a right triangle equals the sum of the squares 
of the lengths of the other two sides.
\end{thm}
\section{Epilogue}
A word or two to conclude, and this even includes some inline 
maths:  $R(x,t)\sim t^{-\beta}g(x/t^\alpha)\exp(-|x|/t^\alpha)$.

\begin{ack}
Partially supported by the Roman Senate.
\end{ack}

% \bibliographystyle{plain}        % Include this if you use bibtex 
% \bibliography{autosam}           % and a bib file to produce the 
%                                  % bibliography (preferred). The
%                                  % correct style is generated by
%                                  % Elsevier at the time of printing.

\appendix
\section{A summary of Latin grammar}
\section{Some Latin vocabulary} 
\end{document}